\section{Заключение}
\label{sec:zaklyuchenie}

Първо, прокурорът е централният орган на досъдебното производство. Той ръководи
разследването и упражнява надзор за законосъобразността му, може да разследва сам и
повдига обвинението, а разследващите органи действат под негово ръководство. Прокурорът
определя и началото, и края на производството.

Второ, самостоятелността на прокурора е ограничена от съдебния контрол. Най-тежките
мерки изискват разрешение на съда по мотивирано искане на прокурора, а прекратяването,
спирането и от 2023~г. отказът да се образува производство за тежко престъпление подлежат
на съдебен контрол. Този модел съответства на стандартите на Съда на ЕС и на ЕСПЧ.

Трето, при киберпрестъпленията значението на прокурора нараства. Компютърните
престъпления вече са тежки, затова при разследването им могат да се използват и
способите, които най-силно засягат правата, но всеки от тях минава през искане на
прокурора. Прокурорът отговаря за бързото запазване на данните, за целостта на иззетите
носители и криптоактиви и за събирането на доказателства в чужбина.

Четвърто, ефективността на тази роля зависи от специализацията и от техническата
подготовка на прокурорите, както и от приемането на националните правила за
прилагане на Регламент (ЕС) 2023/1543. Заслужава обсъждане и изричното уреждане на възможността
прокурорът да разпорежда запазване на компютърни данни още при предварителната проверка.
